\section{Formal System and Research Scope}

\subsection{Object Layer, Theorem Layer, and Semantic Layer}

This paper distinguishes three classes of statements:
\begin{enumerate}
  \item Object layer: superspace embeddings, windows (acceptance domains), symbolic processes, graph structures, and measure-theoretic objects;
  \item Theorem layer: cited or proven mathematical conclusions;
  \item Semantic layer: mappings of statistical quantities at the object layer to physical semantic interpretation, serving only as an interpretation interface without reverse feedback to theorem premises.
\end{enumerate}

\subsection{Strict Meaning of ``Three-Dimensional Irrational Structure Closed in Six Dimensions''}

The online crystallography dictionary defines Fourier module as a $\ZZ$-module of rank $n$ and provides superspace embeddings: if the Fourier module of a three-dimensional quasiperiodic function has rank $n$, then there exist an $n$-dimensional superspace lattice-periodic function $\rho_s$ such that the three-dimensional quasiperiodic function $\rho$ satisfies $\rho(r)=\rho_s(r,0)$, and the superspace dimension $n$ equals the Fourier module rank \cite{IUCRSuperspace}.
Therefore, our claim that the structure ``closes in $6$ dimensions'' strictly restricts to model classes of rank $n=6$; we make no upper bound assertion regarding the minimal superspace dimension for general aperiodic structures.

\subsection{Fixed Bandwidth in Main Text and Parametrization of General $m$}
Unless otherwise stated, the main text fixes the finite window length $m=6$, and relegates parametrization with general $m$ to the appendix or reproducibility protocol sections. The purpose of this convention is to place ``resolution enhancement'' entirely at the protocol level, realized through recursive readout (partition refinement), without depending on increasing $m$, while keeping the geometric model class (rank $6$) unchanged.
Correspondingly, if the notation $x_t^{(6)}$ appears in the main text, it only denotes the output under fixed bandwidth and may be used interchangeably with $x_t$.
